\chapter{Red-Black Trees}

In the previous chapter you met the binary search tree (BST): a delightfully
simple idea where every key sits to the left of bigger keys and to the right of
smaller ones, so searching feels like a game of higher-or-lower. The catch, you
may recall, was painful. A plain BST is only fast \emph{if it stays bushy}. Feed
it keys in sorted order and it quietly degenerates into a linked list, and your
beautiful $\Th(\log n)$ search collapses into a sluggish $\Th(n)$ crawl.

This chapter is about fixing that, permanently, with a structure that
\emph{refuses} to get tall and stringy no matter what you throw at it. The
red-black tree is a binary search tree wearing a small amount of bookkeeping --
one bit of color per node -- and that single bit, governed by five quiet rules,
guarantees the tree's height never exceeds roughly $2\log_2 n$. Every operation
stays $\Th(\log n)$, forever, in the worst case. No randomization, no luck, no
``well, on average.'' Just a guarantee.

It is also one of the first algorithms a beginner meets that feels genuinely
intricate. There are cases. There are sub-cases. There are rotations that swing
whole subtrees around. So we will go slowly, draw a lot of pictures, and make
sure that by the end you do not merely \emph{believe} red-black trees work --
you can \emph{see} why they must.

\section{Why Balance Matters}\label{ch13:sec-balance}

Let us feel the problem from the ground up, since everything in this chapter
exists to escape it.

\begin{intuitionbox}
A binary search tree is a sorted structure you walk \emph{down} to find things.
Each step of the walk should cut the remaining possibilities roughly in half --
that is what makes search logarithmic. But ``cut in half'' only happens if, at
each node, the two subtrees hanging below are about the same size. If one side is
always tiny and the other always huge, you barely make progress per step. The
whole speed of a BST lives or dies on this one question: \emph{is it balanced?}
\end{intuitionbox}

Concretely, suppose you insert the keys $1, 2, 3, 4, 5$ in that order into an
ordinary BST. The first key, $1$, becomes the root. Then $2$ is bigger, so it
goes right. Then $3$ is bigger than both, so it goes right again. And so on. You
get this:

\begin{center}
\begin{tikzpicture}[level distance=10mm,
  every node/.style={circle,draw,minimum size=7mm,inner sep=0pt},
  edge from parent/.style={draw,-{Latex[length=2mm]}}]
\node {1}
  child[missing]
  child {node {2}
    child[missing]
    child {node {3}
      child[missing]
      child {node {4}
        child[missing]
        child {node {5}}}}};
\end{tikzpicture}
\end{center}

That is not a tree in any useful sense; it is a list drawn diagonally. Searching
for $5$ visits all five nodes. With a million sorted keys you would visit a
million nodes. The $\Th(\log n)$ dream is dead.

Now compare it to the \emph{balanced} arrangement of the same five keys:

\begin{center}
\begin{tikzpicture}[level distance=11mm, level 1/.style={sibling distance=24mm},
  level 2/.style={sibling distance=14mm},
  every node/.style={circle,draw,minimum size=7mm,inner sep=0pt},
  edge from parent/.style={draw,-{Latex[length=2mm]}}]
\node {2}
  child {node {1}}
  child {node {4}
    child {node {3}}
    child {node {5}}};
\end{tikzpicture}
\end{center}

Here finding $5$ takes two steps. Finding anything takes at most three. The keys
are identical; only the \emph{shape} differs, and shape is everything.

\begin{focusbox}
The height of a BST is the length of the longest root-to-leaf path, and the cost
of search, insert, and delete is $\Th(\text{height})$. A balanced tree on $n$
keys has height $\Th(\log n)$; a degenerate one has height $\Th(n)$. The entire
goal of this chapter is to keep the height $O(\log n)$ \emph{automatically},
after every single insertion and deletion, without ever rebuilding the tree.
\end{focusbox}

\begin{aestheticbox}
Why should you care beyond passing a course? Because balanced search trees are
\emph{everywhere} the moment a system needs ordered, dynamic data with hard
worst-case guarantees. The C++ \texttt{std::map} and \texttt{std::set}, Java's
\texttt{TreeMap}, the Linux kernel's process scheduler (which once used a
red-black tree to pick the next task to run), the data structures inside
databases that answer range queries, virtual-memory region trackers -- all lean
on red-black trees or close cousins. The reason engineers reach for them over a
hash table is precisely the \emph{ordering}: a hash table can tell you whether
$42$ is present, but a red-black tree can also tell you the smallest key larger
than $42$, walk the keys in sorted order, and do it all with a guaranteed worst
case rather than a hopeful average.
\end{aestheticbox}

The naive fix would be to rebalance from scratch whenever the tree gets lopsided,
but rebuilding costs $\Th(n)$ and we cannot afford that on every insert. We need
\emph{local} surgery: small, $O(1)$ adjustments near the point of change that
keep the whole tree balanced. That surgery has two ingredients -- a set of color
rules that \emph{certify} balance, and an operation called a rotation that lets
us \emph{restore} the rules cheaply. Both are coming up.

\section{The Red-Black Properties}

A red-black tree is a binary search tree in which every node carries one extra
attribute: a color, either red or black. We also adopt one small but important
convention that will save us endless case-fiddling later.

\begin{focusbox}
\textbf{The sentinel.} Instead of having children pointers that are sometimes
\texttt{NIL}, every ``missing'' child and the parent of the root all point to a
single shared dummy node we call \emph{the sentinel}, written $T.nil$. The
sentinel is colored \emph{black} and its key is irrelevant. Treating every leaf
as this one black node means we never have to ask ``but what if this pointer is
null?'' in the middle of a delicate argument. Every real (internal) node thus has
exactly two children, each either a real node or the sentinel.
\end{focusbox}

Now the rules. A binary search tree is a \emph{red-black tree} if it satisfies all
five of the following \emph{red-black properties}.

\begin{formalbox}
\textbf{The five red-black properties.}
\begin{enumerate}
\item Every node is either red or black.
\item The root is black.
\item Every leaf (the sentinel $T.nil$) is black.
\item If a node is red, then both of its children are black. (No red node has a
red child -- equivalently, ``no two reds in a row'' down any path.)
\item For each node, all simple paths from that node down to descendant leaves
contain the same number of black nodes.
\end{enumerate}
\end{formalbox}

Properties 1, 2, and 3 are mostly housekeeping. The two that do the real work are
4 and 5, and they deserve some intuition.

\begin{intuitionbox}
Think of black nodes as the structure's ``load-bearing'' nodes and red nodes as
``free'' decoration squeezed in between. Property 5 says: \emph{measured in black
nodes only, the tree is perfectly balanced} -- every path from a node to a leaf
crosses exactly the same number of black nodes. That is a very strong balance
condition. On its own, though, property 5 would let you pad one path with a long
run of red nodes and make it much longer than another. Property 4 forbids that:
since no two reds may touch, the reds on any path can at most \emph{double} its
length compared to the blacks alone. Put the two together and you get: the
longest path is at most twice the shortest. That ``at most twice'' is the whole
ballgame.
\end{intuitionbox}

This motivates a piece of vocabulary we will use constantly.

\begin{definition}[Black-height]\label{ch13:def-black-height}
The \emph{black-height} of a node $x$, written $\mathrm{bh}(x)$, is the number of
black nodes on any simple path from $x$ down to a leaf, \emph{not counting $x$
itself} but counting the leaf sentinel. Property 5 guarantees this number is the
same no matter which downward path you take, so it is well-defined.
\end{definition}

Here is a small valid red-black tree. Black nodes are drawn as solid dark
circles, red nodes as lighter circles with a red border. The little squares are
the sentinel leaves.

\begin{center}
\begin{tikzpicture}[level distance=12mm,
  level 1/.style={sibling distance=34mm},
  level 2/.style={sibling distance=18mm},
  level 3/.style={sibling distance=10mm},
  black/.style={circle,draw,fill=black,text=white,minimum size=7mm,inner sep=0pt},
  red/.style={circle,draw=trapred,very thick,fill=red!12,minimum size=7mm,inner sep=0pt},
  nil/.style={rectangle,draw,fill=black,text=white,minimum size=4mm,inner sep=1pt,font=\tiny},
  edge from parent/.style={draw,-{Latex[length=1.6mm]}}]
\node[black] {13}
  child {node[red] {8}
    child {node[black] {1}
      child {node[nil] {nil}}
      child {node[red] {6}
        child {node[nil] {nil}}
        child {node[nil] {nil}}}}
    child {node[black] {11}
      child {node[nil] {nil}}
      child {node[nil] {nil}}}}
  child {node[red] {17}
    child {node[black] {15}
      child {node[nil] {nil}}
      child {node[nil] {nil}}}
    child {node[black] {25}
      child {node[nil] {nil}}
      child {node[nil] {nil}}}};
\end{tikzpicture}
\end{center}

Walk through the properties on this tree to convince yourself it is legal. The
root $13$ is black (property 2). No red node has a red child -- $8$ is red and its
children $1$ and $11$ are both black; $17$ is red and its children $15$ and $25$
are both black; $6$ is red and both its children are the black sentinel (property
4). And property 5: pick the root and count black nodes (including sentinels) on
each downward path. Down through $8, 1, 6$ to a sentinel you cross blacks
$\{1, \text{sentinel}\}$ -- wait, let us count carefully, since this is exactly
the kind of counting people botch.

\begin{labbox}
\textbf{Lab: count the black-height by hand.} Take the root $13$. Trace the path
$13 \to 8 \to 1 \to 6 \to \text{nil}$. The nodes \emph{below} $13$ on this path
are $8$ (red), $1$ (black), $6$ (red), and the sentinel (black). Black nodes,
not counting $13$ itself: $1$ and the sentinel, so $2$. Now trace
$13 \to 17 \to 15 \to \text{nil}$: below $13$ we have $17$ (red), $15$ (black),
sentinel (black) -- again $2$ blacks. Try every downward path from $13$ you can
find; you will always get $2$. So $\mathrm{bh}(13) = 2$. Do the same starting
from node $8$: every path below it hits exactly $1$ black plus its sentinel,
giving $\mathrm{bh}(8) = 2$ as well (paths through $1$ cross $1$ then sentinel;
path to $11$ crosses $11$ then sentinel). Counting the sentinels is what keeps
property 5 working out -- skip them and you will miscount.
\end{labbox}

\begin{trapbox}
Three beginner snares with the properties:
\begin{itemize}
\item \textbf{Forgetting the sentinel counts.} Property 5 counts the black leaf at
the bottom of every path. Many off-by-one errors come from omitting it.
\item \textbf{Thinking ``balanced'' means equal subtree sizes.} It does not.
Red-black balance is about \emph{black-height equality}, not node counts. The
left subtree can hold more nodes than the right; that is fine.
\item \textbf{Reading property 4 as ``red children must be black.''} It is the
other direction: a \emph{red parent} forces \emph{black children}. A black node is
free to have children of either color.
\end{itemize}
\end{trapbox}

\section{The Height Bound}

Before we touch insertion or deletion, let us cash in the properties and prove
the payoff: a red-black tree on $n$ internal nodes has height at most
$2\log_2(n+1)$. This is the theorem that makes the whole structure worth its
trouble, and the proof is a short induction.

\begin{lemma}[Subtree size from black-height]\label{ch13:lem-bh-size}
The subtree rooted at any node $x$ contains at least $2^{\mathrm{bh}(x)} - 1$
internal nodes.
\end{lemma}

\begin{intuitionbox}
The claim says: ``a node with black-height $b$ is sitting on top of at least a
full black tree of height $b$.'' A perfectly full binary tree of height $b$ has
$2^b - 1$ internal nodes, and red nodes can only \emph{add} to that count, never
subtract. So the black-height alone forces a minimum population. The bigger the
black-height, the more nodes you are guaranteed to have crammed underneath.
\end{intuitionbox}

\begin{proof}
We induct on the height of $x$ (the ordinary height, longest path down to a
leaf).

\emph{Base case.} If $x$ has height $0$, then $x$ is the sentinel itself. Its
black-height is $0$, and the formula demands at least $2^0 - 1 = 0$ internal
nodes. The sentinel's subtree has zero internal nodes. True.

\emph{Inductive step.} Let $x$ be an internal node of positive height with two
children. Each child has height strictly less than $x$, so the induction
hypothesis applies to both. Now, what is each child's black-height? A child is
either red or black. If the child is black it contributes $1$ to the black count
on paths through it, so its own black-height is $\mathrm{bh}(x) - 1$. If the child
is red it contributes $0$, so its black-height is $\mathrm{bh}(x)$. Either way,
each child has black-height \emph{at least} $\mathrm{bh}(x) - 1$. By the induction
hypothesis, each child's subtree holds at least $2^{\mathrm{bh}(x)-1} - 1$
internal nodes. Adding the two children's subtrees plus $x$ itself,
\[
  \big(2^{\mathrm{bh}(x)-1} - 1\big) + \big(2^{\mathrm{bh}(x)-1} - 1\big) + 1
  = 2 \cdot 2^{\mathrm{bh}(x)-1} - 1
  = 2^{\mathrm{bh}(x)} - 1,
\]
which is exactly what we wanted to show.
\end{proof}

\begin{theorem}[Height of a red-black tree]\label{ch13:thm-height}
A red-black tree with $n$ internal nodes has height at most $2\log_2(n+1)$.
\end{theorem}

\begin{proof}
Let $h$ be the height of the tree, that is, the length of the longest path from
the root down to a leaf. Property 4 (no two reds in a row) means that at least
half the nodes on that longest path are black: as you walk down, you cannot take
two red steps back to back, so reds are at most every other node. Hence the
black-height of the root is at least $h/2$:
\[
  \mathrm{bh}(\text{root}) \ge \frac{h}{2}.
\]
By Lemma~\ref{ch13:lem-bh-size} applied at the root, the whole tree has at least
$2^{\mathrm{bh}(\text{root})} - 1 \ge 2^{h/2} - 1$ internal nodes. So
\[
  n \ge 2^{h/2} - 1 \quad\Longrightarrow\quad n + 1 \ge 2^{h/2}
  \quad\Longrightarrow\quad \log_2(n+1) \ge \frac{h}{2}
  \quad\Longrightarrow\quad h \le 2\log_2(n+1). \qedhere
\]
\end{proof}

The picture below makes the ``black-height pins the height'' idea concrete. Two
paths leave the same root. The left one is the densest a path can be: black and
red strictly alternating, which property 4 permits but cannot exceed. The right
one is the sparsest: all black. Both must carry the \emph{same} number of black
nodes (property 5), so the alternating path can be at most twice as long.

\begin{center}
\begin{tikzpicture}[
  black/.style={circle,draw,fill=black,text=white,minimum size=5.5mm,inner sep=0pt,font=\scriptsize},
  red/.style={circle,draw=trapred,very thick,fill=red!12,minimum size=5.5mm,inner sep=0pt,font=\scriptsize},
  every edge/.style={draw,-{Latex[length=1.4mm]}}]
\node[black] (r) at (0,0) {};
% long alternating path on the left
\node[red]   (a1) at (-1.6,-1.1) {};
\node[black] (a2) at (-1.6,-2.2) {};
\node[red]   (a3) at (-1.6,-3.3) {};
\node[black] (a4) at (-1.6,-4.4) {};
\draw (r) edge (a1); \draw (a1) edge (a2); \draw (a2) edge (a3); \draw (a3) edge (a4);
% short all-black path on the right
\node[black] (b1) at (1.6,-1.1) {};
\node[black] (b2) at (1.6,-2.2) {};
\draw (r) edge (b1); \draw (b1) edge (b2);
% brackets / labels
\node[trapred,font=\scriptsize,align=center] at (-3.5,-2.2)
  {longest path\\ (red/black\\ alternating)};
\node[font=\scriptsize,align=center] at (3.4,-1.65)
  {shortest path\\ (all black)};
\node[font=\scriptsize,align=center] at (0,-5.2)
  {both paths cross the same number of \emph{black} nodes (property 5);\\
   the alternating path is at most twice as long (property 4).};
\end{tikzpicture}
\end{center}

\begin{focusbox}
That is the prize. Because the height is $O(\log n)$, every operation that runs in
time proportional to the height -- search, minimum, maximum, successor,
predecessor, insert, delete -- runs in $O(\log n)$ \emph{worst case}. Searching a
red-black tree is identical to searching an ordinary BST (the colors are
invisible to a search); it just so happens that the tree you are searching is
guaranteed to be short.
\end{focusbox}

\begin{walkbox}
Here is something to chew on while you stretch your legs. The bound says the
longest path is at most twice the shortest. Could you build a \emph{valid}
red-black tree where the longest root-to-leaf path is exactly twice the shortest?
What would it have to look like -- where do the red nodes go, and where are they
forbidden? Try to sketch the smallest such tree. (Hint: alternate red and black
down the long path, and keep the short path all black.)
\end{walkbox}

\section{Rotations}\label{ch13:sec-rotations}

Searching never changes the tree, so it never breaks a property. Insertion and
deletion \emph{do} change the tree, and they will temporarily break the rules. To
repair them we need a tool that restructures the tree locally while preserving the
binary-search-tree ordering. That tool is the \emph{rotation}.

\begin{intuitionbox}
A rotation grabs an edge between a parent and a child and ``pivots'' it, swapping
which of the two is on top -- like rotating a seesaw. The magic is that it can be
done while keeping all the keys in valid search-tree order, because there is
exactly one subtree caught in the middle that has to change hands, and it lands in
a spot that respects the ordering. A left rotation makes the right child rise; a
right rotation makes the left child rise. Each touches only a constant number of
pointers, so it costs $O(1)$.
\end{intuitionbox}

Picture a left rotation around node $x$, whose right child is $y$. We want $y$ to
rise above $x$. The subtree $\beta$ (which sits between $x$ and $y$ in key order)
is the one that has to move.

\begin{center}
\begin{tikzpicture}[
  node/.style={circle,draw,minimum size=8mm,inner sep=0pt},
  tri/.style={isosceles triangle,draw,shape border rotate=90,minimum width=10mm,minimum height=7mm,inner sep=1pt,font=\small},
  every edge/.style={draw,-{Latex[length=1.6mm]}}]
% left configuration
\node[node] (x) at (0,0) {$x$};
\node[node] (y) at (1.5,-1.6) {$y$};
\node[tri] (a) at (-1.4,-1.6) {$\alpha$};
\node[tri] (b) at (0.4,-3.0) {$\beta$};
\node[tri] (g) at (2.6,-3.0) {$\gamma$};
\draw (x) edge (a);
\draw (x) edge (y);
\draw (y) edge (b);
\draw (y) edge (g);
% arrow
\node at (4.4,-1.5) {\Large $\xrightarrow{\;\text{left-rotate}(x)\;}$};
% right configuration
\node[node] (y2) at (8.4,0) {$y$};
\node[node] (x2) at (6.9,-1.6) {$x$};
\node[tri] (g2) at (9.8,-1.6) {$\gamma$};
\node[tri] (a2) at (5.8,-3.0) {$\alpha$};
\node[tri] (b2) at (7.9,-3.0) {$\beta$};
\draw (y2) edge (x2);
\draw (y2) edge (g2);
\draw (x2) edge (a2);
\draw (x2) edge (b2);
\end{tikzpicture}
\end{center}

Read the key ordering off both pictures. Before: $\alpha < x < \beta < y <
\gamma$. After: $\alpha < x < \beta < y < \gamma$. \emph{Identical.} The rotation
slid $\beta$ from being $y$'s left child to being $x$'s right child, which is
exactly where it still belongs, because everything in $\beta$ is greater than $x$
and less than $y$ in both pictures. The in-order traversal is unchanged; only the
shape moved.

A right rotation is the mirror image: it is the inverse operation, taking the
right picture back to the left. Here is the code for a left rotation, written
against the sentinel convention so we never test for null.

\begin{algorithm}[H]
\caption{\textsc{Left-Rotate}$(T, x)$}
\KwIn{A red-black tree $T$ and a node $x$ whose right child is not the sentinel.}
\KwOut{The tree, restructured so $x$'s former right child rises above it.}
$y \gets x.\mathit{right}$\;
$x.\mathit{right} \gets y.\mathit{left}$ \tcp*{$\beta$ becomes $x$'s right subtree}
\If{$y.\mathit{left} \neq T.\mathit{nil}$}{$y.\mathit{left}.\mathit{parent} \gets x$\;}
$y.\mathit{parent} \gets x.\mathit{parent}$ \tcp*{link $y$ to $x$'s old parent}
\uIf{$x.\mathit{parent} = T.\mathit{nil}$}{$T.\mathit{root} \gets y$\;}
\uElseIf{$x = x.\mathit{parent}.\mathit{left}$}{$x.\mathit{parent}.\mathit{left} \gets y$\;}
\Else{$x.\mathit{parent}.\mathit{right} \gets y$\;}
$y.\mathit{left} \gets x$ \tcp*{put $x$ on $y$'s left}
$x.\mathit{parent} \gets y$\;
\end{algorithm}

\textsc{Right-Rotate} is obtained by swapping every \emph{left} with \emph{right}
in the code above. Notice there are no loops: a fixed handful of pointer writes,
so each rotation is firmly $O(1)$.

\begin{labbox}
\textbf{Lab: rotate by hand.} Take this little tree (all colors ignored for the
moment; we are just practicing the mechanics):
\begin{center}
\begin{tikzpicture}[level distance=10mm, level 1/.style={sibling distance=20mm},
  level 2/.style={sibling distance=12mm},
  every node/.style={circle,draw,minimum size=7mm,inner sep=0pt},
  edge from parent/.style={draw,-{Latex[length=1.6mm]}}]
\node {7}
  child {node {4}}
  child {node {11}
    child {node {9}}
    child {node {18}}};
\end{tikzpicture}
\end{center}
Left-rotate around the root $7$. Here $x = 7$, $y = 11$, and the subtree $\beta$ is
the node $9$ (it lies between $7$ and $11$). After the rotation, $11$ becomes the
root, $7$ becomes its left child, $9$ moves to be the \emph{right} child of $7$,
and $18$ stays as $11$'s right child. Draw it and read the in-order traversal of
both: you should get $4, 7, 9, 11, 18$ both times. If your before-and-after
in-orders disagree, you moved the wrong subtree.
\end{labbox}

\begin{trapbox}
The single most common rotation bug is forgetting to reattach \emph{three}
parent pointers: the moved subtree's parent ($\beta$ now points up to $x$), $y$'s
parent (now $x$'s old parent), and $x$'s parent (now $y$). Miss any one and the
tree's downward pointers will look fine while the upward ones lie -- and the
fixup code that walks \emph{up} the tree will wander off into the weeds.
\end{trapbox}

\section{Insertion}

We now have everything we need. Insertion happens in two acts. In Act One we drop
the new key in exactly where an ordinary BST would put it, and we paint it
\emph{red}. In Act Two we walk back up the tree fixing any property the red paint
job may have violated.

\begin{intuitionbox}
Why color the new node red rather than black? Because red is the ``cheap'' color:
a red node does not change anybody's black-height. If we inserted a black node, we
would instantly break property 5 along that one path (it now has one more black
than its siblings) and the repair would be ugly. By inserting red, the only
property we can possibly violate is property 4 (a red node with a red parent) --
and property 2 if the new node happens to be the root. Those are far gentler to
fix. So we always start red and only darken nodes as the fixup demands.
\end{intuitionbox}

The first act is just BST insertion with the sentinel bookkeeping; here it is.

\begin{algorithm}[H]
\caption{\textsc{RB-Insert}$(T, z)$}
\KwIn{Tree $T$ and a fresh node $z$ holding the key to insert.}
\KwOut{$T$ with $z$ inserted and all five properties restored.}
$y \gets T.\mathit{nil}$; \quad $x \gets T.\mathit{root}$\;
\While{$x \neq T.\mathit{nil}$}{
  $y \gets x$\;
  \uIf{$z.\mathit{key} < x.\mathit{key}$}{$x \gets x.\mathit{left}$\;}
  \Else{$x \gets x.\mathit{right}$\;}
}
$z.\mathit{parent} \gets y$\;
\uIf{$y = T.\mathit{nil}$}{$T.\mathit{root} \gets z$\;}
\uElseIf{$z.\mathit{key} < y.\mathit{key}$}{$y.\mathit{left} \gets z$\;}
\Else{$y.\mathit{right} \gets z$\;}
$z.\mathit{left} \gets T.\mathit{nil}$; \quad $z.\mathit{right} \gets T.\mathit{nil}$; \quad $z.\mathit{color} \gets \textsc{red}$\;
\textsc{RB-Insert-Fixup}$(T, z)$\;
\end{algorithm}

Everything interesting is in the fixup, which involves a few nodes near $z$: its
parent, and its parent's sibling, which we call the \emph{uncle}. The uncle's
color is the hinge on which the whole case analysis turns.

\begin{focusbox}
After Act One, the only thing that can be wrong is a \emph{red-red} violation:
$z$ is red, and its parent might also be red. (If the parent is black we are
already done; if $z$ is the root we just paint it black.) So the fixup loop's job
is: ``as long as $z$ has a red parent, push the violation upward or rotate it
away.'' Because each pass either moves the problem two levels closer to the root
or ends with a rotation, the loop runs $O(\log n)$ times.
\end{focusbox}

There are three cases, and they come in a left-leaning flavor and a mirror-image
right-leaning flavor. We describe the version where $z$'s parent is a \emph{left}
child of the grandparent; the other side is symmetric (swap left and right
throughout). Let $z.p$ be the parent, $z.p.p$ the grandparent, and $y$ the uncle.

\textbf{Case 1: the uncle $y$ is red.} Both the parent and the uncle are red,
sitting under a (necessarily black) grandparent. The fix is pure recoloring, no
rotation: paint the parent and the uncle black, paint the grandparent red, then
move our attention up to the grandparent and continue. We have ``pushed'' the
red-red conflict up two levels.

\begin{center}
\begin{tikzpicture}[
  black/.style={circle,draw,fill=black,text=white,minimum size=7mm,inner sep=0pt},
  red/.style={circle,draw=trapred,very thick,fill=red!12,minimum size=7mm,inner sep=0pt},
  tri/.style={isosceles triangle,draw,shape border rotate=90,minimum width=7mm,minimum height=5mm,inner sep=0.5pt,font=\scriptsize},
  every edge/.style={draw,-{Latex[length=1.4mm]}}]
% before
\node[black] (gp) at (0,0) {$C$};
\node[red] (p) at (-1.2,-1.4) {$A$};
\node[red] (u) at (1.2,-1.4) {$D$};
\node[red] (z) at (-2.0,-2.8) {$z$};
\draw (gp) edge (p); \draw (gp) edge (u); \draw (p) edge (z);
\node at (3.0,-1.4) {\large $\Rightarrow$};
% after
\node[red] (gp2) at (6,0) {$C$};
\node[black] (p2) at (4.8,-1.4) {$A$};
\node[black] (u2) at (7.2,-1.4) {$D$};
\node[red] (z2) at (4.0,-2.8) {$z$};
\draw (gp2) edge (p2); \draw (gp2) edge (u2); \draw (p2) edge (z2);
\node at (6,-3.6) {\scriptsize $z$ moves up to $C$};
\end{tikzpicture}
\end{center}

Why does this preserve property 5? The grandparent went from black to red and its
two children went from red to black. Any path passing through this little cluster
used to cross the black grandparent and then one of the now-black children -- the
total black count along every such path is unchanged, because we removed one black
(the grandparent) and added one back (a child). The black-height is conserved
exactly. We may have created a new red-red conflict higher up (if the grandparent's
own parent is red), which is why we loop with $z$ now pointing at the grandparent.

\textbf{Case 2: the uncle $y$ is black, and $z$ is a ``zig-zag'' (an inner
grandchild).} That is, $z$ is the right child of a left-child parent. A single
rotation will not straighten this; we first do a \emph{left rotation on the
parent} to turn the zig-zag into a straight line, converting it into Case 3.
After this rotation we relabel so that $z$ points at what is now the lower node.

\textbf{Case 3: the uncle $y$ is black, and $z$ is a straight line (an outer
grandchild).} Now $z$ is the left child of a left-child parent -- parent and
grandparent and $z$ all lean the same way. Recolor the parent black and the
grandparent red, then \emph{right-rotate the grandparent}. This lifts the parent
into the grandparent's old slot, with the (formerly red, now black) parent on top
and the grandparent and $z$ as its red children -- no two reds touch, and the
loop terminates.

\begin{center}
\begin{tikzpicture}[
  black/.style={circle,draw,fill=black,text=white,minimum size=7mm,inner sep=0pt},
  red/.style={circle,draw=trapred,very thick,fill=red!12,minimum size=7mm,inner sep=0pt},
  every edge/.style={draw,-{Latex[length=1.4mm]}}]
% case 2 -> case 3 -> done
\node[black] (gp) at (0,0) {$C$};
\node[red] (p) at (-1.1,-1.3) {$A$};
\node[red] (z) at (-0.1,-2.6) {$z{=}B$};
\draw (gp) edge (p); \draw (p) edge (z);
\node at (2.0,-1.3) {\large $\xrightarrow{\text{LR}(A)}$};
% case 3 config
\node[black] (gp2) at (4.4,0) {$C$};
\node[red] (b) at (3.3,-1.3) {$B$};
\node[red] (a) at (2.3,-2.6) {$A$};
\draw (gp2) edge (b); \draw (b) edge (a);
\node at (6.0,-1.3) {\large $\xrightarrow{\text{RR}(C)}$};
% done
\node[black] (b2) at (8.4,-0.6) {$B$};
\node[red] (a2) at (7.4,-2.0) {$A$};
\node[red] (c2) at (9.4,-2.0) {$C$};
\draw (b2) edge (a2); \draw (b2) edge (c2);
\end{tikzpicture}
\end{center}

Here is the whole fixup, both sides.

\begin{algorithm}[H]
\caption{\textsc{RB-Insert-Fixup}$(T, z)$}
\KwIn{Tree $T$ where the only possible violation is a red-red edge at $z$.}
\KwOut{$T$ with all five properties restored.}
\While{$z.\mathit{parent}.\mathit{color} = \textsc{red}$}{
  \uIf{$z.\mathit{parent} = z.\mathit{parent}.\mathit{parent}.\mathit{left}$}{
    $y \gets z.\mathit{parent}.\mathit{parent}.\mathit{right}$ \tcp*{the uncle}
    \uIf(\tcp*[f]{Case 1}){$y.\mathit{color} = \textsc{red}$}{
      $z.\mathit{parent}.\mathit{color} \gets \textsc{black}$\;
      $y.\mathit{color} \gets \textsc{black}$\;
      $z.\mathit{parent}.\mathit{parent}.\mathit{color} \gets \textsc{red}$\;
      $z \gets z.\mathit{parent}.\mathit{parent}$\;
    }\Else{
      \If(\tcp*[f]{Case 2}){$z = z.\mathit{parent}.\mathit{right}$}{
        $z \gets z.\mathit{parent}$\;
        \textsc{Left-Rotate}$(T, z)$\;
      }
      $z.\mathit{parent}.\mathit{color} \gets \textsc{black}$ \tcp*{Case 3}
      $z.\mathit{parent}.\mathit{parent}.\mathit{color} \gets \textsc{red}$\;
      \textsc{Right-Rotate}$(T, z.\mathit{parent}.\mathit{parent})$\;
    }
  }\Else{
    \emph{same as above with \textsc{left} and \textsc{right} exchanged}\;
  }
}
$T.\mathit{root}.\mathit{color} \gets \textsc{black}$\;
\end{algorithm}

\begin{focusbox}
The very last line -- forcibly painting the root black -- quietly handles two
situations at once: the new node was inserted into an empty tree and became the
root, or Case 1 propagated a red all the way up to the root. Painting the root
black can never violate property 5 (it adds one black to \emph{every} path
equally) and it restores property 2. It is the tidy bow on top.
\end{focusbox}

\subsection{A fully worked insertion}\label{ch13:sec-worked-insert}

We will build a red-black tree from empty by inserting $10, 20, 30, 15, 25, 5, 35$
and watch the fixups fire, narrating the interesting moments.

\begin{enumerate}
\item \textbf{Insert $10$.} Empty tree, so $10$ becomes the root, inserted red,
then the final line paints it black. Tree: black $10$.
\item \textbf{Insert $20$.} Goes right of $10$, red. Its parent $10$ is black, no
violation. Tree: black $10$ with red right child $20$.
\item \textbf{Insert $30$.} Goes right of $20$, red. Now $20$ (red) has a red
child $30$: a red-red violation, and $20$'s uncle is the sentinel (black). $z=30$
is an outer (right-right) grandchild, so this is Case 3 on the right side:
recolor $20$ black, $10$ red, and left-rotate the root $10$. Now $20$ is the black
root with red children $10$ and $30$. Balanced and legal.
\item \textbf{Insert $15$.} From $20$ go left to $10$, then right of $10$, red.
Parent $10$ is red -- violation. The uncle here is $30$, which is \emph{red}.
Case 1: paint $10$ and $30$ black, paint $20$ red, move $z$ up to $20$. But $20$
is the root, and its parent (the sentinel) is black, so the loop stops. The final
line paints the root $20$ black again. Tree now: black $20$; black $10$ (with red
right child $15$); black $30$.
\item \textbf{Insert $25$.} From $20$ right to $30$, then left of $30$, red.
Parent $30$ is black -- no violation. Easy.
\item \textbf{Insert $5$.} From $20$ left to $10$, left of $10$, red. Parent $10$
is black -- no violation.
\item \textbf{Insert $35$.} From $20$ right to $30$, right of $30$, red. Parent
$30$ is black -- no violation.
\end{enumerate}

The final tree:

\begin{center}
\begin{tikzpicture}[level distance=12mm,
  level 1/.style={sibling distance=34mm},
  level 2/.style={sibling distance=17mm},
  black/.style={circle,draw,fill=black,text=white,minimum size=7mm,inner sep=0pt},
  red/.style={circle,draw=trapred,very thick,fill=red!12,minimum size=7mm,inner sep=0pt},
  edge from parent/.style={draw,-{Latex[length=1.6mm]}}]
\node[black] {20}
  child {node[black] {10}
    child {node[red] {5}}
    child {node[red] {15}}}
  child {node[black] {30}
    child {node[red] {25}}
    child {node[red] {35}}};
\end{tikzpicture}
\end{center}

Notice how the sorted-order insertion that destroyed a plain BST in
Section~\ref{ch13:sec-balance} would, here, have been quietly rebalanced at every step. That is the
whole point.

\begin{complexitybox}
\textbf{Insertion cost.} Act One (the BST descent) is $O(\log n)$ because the
tree's height is $O(\log n)$. The fixup loop runs $O(\log n)$ iterations in the
worst case (Case 1 climbs two levels per iteration), and crucially it performs
\emph{at most two rotations} total -- Cases 2 and 3 each end the loop. So insertion
is $O(\log n)$ time with $O(1)$ rotations. The low rotation count is what makes
red-black trees attractive in practice: restructuring is cheap.
\end{complexitybox}

\section{Deletion}

Deletion is the harder sibling. Searching never breaks anything and insertion can
only create a too-light splash of red; but deletion can \emph{remove} a black node,
and removing a black node punches a hole in property 5 -- suddenly one path has
fewer black nodes than its peers. Repairing that hole is the delicate part.

\begin{intuitionbox}
Removing a \emph{red} node is harmless: red nodes do not count toward
black-height, so plucking one out leaves every path's black count unchanged, and
no red-red pair can be newly created in the spot it vacated. The trouble is
removing a \emph{black} node. Imagine the missing black node leaves behind a
``debt'' of one black on its path. We model this with a colorful fiction: the node
that moves into the gap is said to carry an \emph{extra black} -- it becomes
``doubly black'' -- and the entire fixup is the process of \emph{discharging that
extra black}, either by shoving it onto a red node nearby (which can simply absorb
it by turning black) or by sliding it upward until it reaches the root or a red
node that can swallow it.
\end{intuitionbox}

As with ordinary BST deletion, the node we physically remove is not always the
node holding the key we asked to delete. If the target has two children, we splice
out its \emph{successor} (the smallest key in its right subtree, which has at most
one child) and move that key up. So we track two things: the node $y$ actually
spliced out, and the node $x$ that moves into $y$'s position. The rule is:
\emph{the fixup is needed exactly when the spliced-out node $y$ was black.}

We will lean on a helper, \textsc{RB-Transplant}, that replaces one subtree with
another by rewiring the parent's pointer -- the same plumbing used in plain BST
deletion, adjusted for the sentinel.

\begin{algorithm}[H]
\caption{\textsc{RB-Transplant}$(T, u, v)$}
\KwIn{Replace the subtree rooted at $u$ with the subtree rooted at $v$.}
\uIf{$u.\mathit{parent} = T.\mathit{nil}$}{$T.\mathit{root} \gets v$\;}
\uElseIf{$u = u.\mathit{parent}.\mathit{left}$}{$u.\mathit{parent}.\mathit{left} \gets v$\;}
\Else{$u.\mathit{parent}.\mathit{right} \gets v$\;}
$v.\mathit{parent} \gets u.\mathit{parent}$\;
\end{algorithm}

\begin{algorithm}[H]
\caption{\textsc{RB-Delete}$(T, z)$}
\KwIn{Tree $T$ and the node $z$ to delete.}
\KwOut{$T$ with $z$ removed and all five properties restored.}
$y \gets z$; \quad $y\text{-}orig\text{-}color \gets y.\mathit{color}$\;
\uIf{$z.\mathit{left} = T.\mathit{nil}$}{
  $x \gets z.\mathit{right}$; \quad \textsc{RB-Transplant}$(T, z, z.\mathit{right})$\;
}\uElseIf{$z.\mathit{right} = T.\mathit{nil}$}{
  $x \gets z.\mathit{left}$; \quad \textsc{RB-Transplant}$(T, z, z.\mathit{left})$\;
}\Else{
  $y \gets \textsc{Tree-Minimum}(z.\mathit{right})$ \tcp*{the successor}
  $y\text{-}orig\text{-}color \gets y.\mathit{color}$\;
  $x \gets y.\mathit{right}$\;
  \uIf{$y.\mathit{parent} = z$}{$x.\mathit{parent} \gets y$\;}
  \Else{
    \textsc{RB-Transplant}$(T, y, y.\mathit{right})$\;
    $y.\mathit{right} \gets z.\mathit{right}$; \quad $y.\mathit{right}.\mathit{parent} \gets y$\;
  }
  \textsc{RB-Transplant}$(T, z, y)$\;
  $y.\mathit{left} \gets z.\mathit{left}$; \quad $y.\mathit{left}.\mathit{parent} \gets y$; \quad $y.\mathit{color} \gets z.\mathit{color}$\;
}
\If{$y\text{-}orig\text{-}color = \textsc{black}$}{\textsc{RB-Delete-Fixup}$(T, x)$\;}
\end{algorithm}

\begin{focusbox}
The key bookkeeping: $y$ is the node that physically leaves the tree's structure
(either $z$ itself, or $z$'s successor when $z$ has two children). $x$ is the node
that moves into $y$'s old spot. We remember $y$'s \emph{original} color. If it was
red, no property broke and we are done. If it was black, $x$ inherits an
\emph{extra black} and we call the fixup to discharge it.
\end{focusbox}

Now the fixup. We describe the case where $x$ is a \emph{left} child; the
right-child cases are mirror images. Let $w$ be $x$'s sibling. Because $x$ carries
an extra black (its path is one black short), $w$'s subtree must be ``black-heavy''
relative to $x$'s -- and that asymmetry is what we exploit. There are four cases.

\textbf{Case 1: the sibling $w$ is red.} We cannot reason cleanly while $w$ is
red, so we first convert it away. Recolor $w$ black and $x$'s parent red, then
left-rotate the parent. Now $x$ has a new, black sibling, and we have reduced to
one of the next three cases.

\textbf{Case 2: $w$ is black and both of $w$'s children are black.} Here neither
side of $w$ can absorb the extra black by rotation, so we ``move the debt up.''
Recolor $w$ red -- this lightens $w$'s side by one black, rebalancing $w$'s
subtree against $x$'s -- and pass the extra black up to the parent by setting
$x \gets x.\mathit{parent}$. The loop continues from there. If we ever reach a red
node (or the root), it absorbs the extra black and we stop.

\begin{center}
\begin{tikzpicture}[
  black/.style={circle,draw,fill=black,text=white,minimum size=6mm,inner sep=0pt,font=\small},
  red/.style={circle,draw=trapred,very thick,fill=red!12,minimum size=6mm,inner sep=0pt,font=\small},
  dbl/.style={circle,draw=black,double,double distance=1pt,fill=black!70,text=white,minimum size=6mm,inner sep=0pt,font=\small},
  every edge/.style={draw,-{Latex[length=1.4mm]}}]
% before: B is parent (unknown color, shown grey-ish via black outline label)
\node[black] (b) at (0,0) {$B$};
\node[dbl]   (x) at (-1.1,-1.4) {$x$};
\node[black] (w) at (1.1,-1.4) {$w$};
\node[black] (c) at (0.3,-2.7) {};
\node[black] (e) at (1.9,-2.7) {};
\draw (b) edge (x); \draw (b) edge (w); \draw (w) edge (c); \draw (w) edge (e);
\node at (3.2,-1.4) {\large $\Rightarrow$};
% after: w turns red, extra black floats up to B (now carries the debt)
\node[dbl]   (b2) at (6.2,0) {$B$};
\node[black] (x2) at (5.1,-1.4) {$x'$};
\node[red]   (w2) at (7.3,-1.4) {$w$};
\node[black] (c2) at (6.5,-2.7) {};
\node[black] (e2) at (8.1,-2.7) {};
\draw (b2) edge (x2); \draw (b2) edge (w2); \draw (w2) edge (c2); \draw (w2) edge (e2);
\node[font=\scriptsize,align=center] at (3.1,-3.4)
  {recolor $w$ red; the extra black rises to $B$,\\ which becomes the new $x$.};
\end{tikzpicture}
\end{center}

\textbf{Case 3: $w$ is black, $w$'s left child is red, and $w$'s right child is
black.} This is the ``wrong-side red'' situation; we rotate it into Case 4.
Recolor $w$'s left child black and $w$ red, then right-rotate $w$. Now $x$'s new
sibling is black with a red \emph{right} child, which is Case 4.

\textbf{Case 4: $w$ is black and $w$'s right child is red.} This is the terminal
case that actually discharges the debt. Color $w$ with the parent's color, color
the parent black, color $w$'s red right child black, and left-rotate the parent.
This single rotation restores the missing black on $x$'s path without disturbing
any other path's count, so we set $x \gets T.\mathit{root}$ to end the loop.

\begin{center}
\begin{tikzpicture}[
  black/.style={circle,draw,fill=black,text=white,minimum size=6.5mm,inner sep=0pt,font=\small},
  red/.style={circle,draw=trapred,very thick,fill=red!12,minimum size=6.5mm,inner sep=0pt,font=\small},
  dbl/.style={circle,draw=black,double,double distance=1pt,fill=black!70,text=white,minimum size=6.5mm,inner sep=0pt,font=\small},
  every edge/.style={draw,-{Latex[length=1.4mm]}}]
% case 4 before
\node[black] (b) at (0,0) {$B$};
\node[dbl] (x) at (-1.1,-1.4) {$x$};
\node[black] (d) at (1.1,-1.4) {$D$};
\node[black] (c) at (0.2,-2.7) {$C$};
\node[red] (e) at (2.0,-2.7) {$E$};
\draw (b) edge (x); \draw (b) edge (d); \draw (d) edge (c); \draw (d) edge (e);
\node at (3.4,-1.4) {\large $\Rightarrow$};
% after
\node[black] (d2) at (6.4,0) {$D$};
\node[black] (b2) at (5.2,-1.4) {$B$};
\node[black] (e2) at (7.6,-1.4) {$E$};
\node[black] (x2) at (4.3,-2.7) {$x$};
\node[black] (c2) at (6.1,-2.7) {$C$};
\draw (d2) edge (b2); \draw (d2) edge (e2); \draw (b2) edge (x2); \draw (b2) edge (c2);
\end{tikzpicture}
\end{center}

The doubly-black $x$ on the left has shed its extra black on the right; the red
node $E$ paid the debt by turning black, and the rotation rebalanced the
black-heights.

\begin{algorithm}[H]
\caption{\textsc{RB-Delete-Fixup}$(T, x)$}
\KwIn{Tree $T$ where $x$ carries an extra black.}
\KwOut{$T$ with the extra black discharged and all properties restored.}
\While{$x \neq T.\mathit{root}$ \textbf{and} $x.\mathit{color} = \textsc{black}$}{
  \uIf{$x = x.\mathit{parent}.\mathit{left}$}{
    $w \gets x.\mathit{parent}.\mathit{right}$\;
    \If(\tcp*[f]{Case 1}){$w.\mathit{color} = \textsc{red}$}{
      $w.\mathit{color} \gets \textsc{black}$; \quad $x.\mathit{parent}.\mathit{color} \gets \textsc{red}$\;
      \textsc{Left-Rotate}$(T, x.\mathit{parent})$; \quad $w \gets x.\mathit{parent}.\mathit{right}$\;
    }
    \uIf(\tcp*[f]{Case 2}){$w.\mathit{left}.\mathit{color} = \textsc{black}$ \textbf{and} $w.\mathit{right}.\mathit{color} = \textsc{black}$}{
      $w.\mathit{color} \gets \textsc{red}$; \quad $x \gets x.\mathit{parent}$\;
    }\Else{
      \If(\tcp*[f]{Case 3}){$w.\mathit{right}.\mathit{color} = \textsc{black}$}{
        $w.\mathit{left}.\mathit{color} \gets \textsc{black}$; \quad $w.\mathit{color} \gets \textsc{red}$\;
        \textsc{Right-Rotate}$(T, w)$; \quad $w \gets x.\mathit{parent}.\mathit{right}$\;
      }
      $w.\mathit{color} \gets x.\mathit{parent}.\mathit{color}$ \tcp*{Case 4}
      $x.\mathit{parent}.\mathit{color} \gets \textsc{black}$; \quad $w.\mathit{right}.\mathit{color} \gets \textsc{black}$\;
      \textsc{Left-Rotate}$(T, x.\mathit{parent})$; \quad $x \gets T.\mathit{root}$\;
    }
  }\Else{
    \emph{same as above with \textsc{left} and \textsc{right} exchanged}\;
  }
}
$x.\mathit{color} \gets \textsc{black}$\;
\end{algorithm}

\begin{focusbox}
A map of the four cases, so they stop feeling arbitrary:
\begin{itemize}
\item \textbf{Case 1} is a \emph{setup} move: it gets rid of a red sibling so the
others can reason about a black one. It always falls through into 2, 3, or 4.
\item \textbf{Case 2} is the only case that \emph{loops}. It cannot fix the debt
locally (both nephews are black), so it pushes the extra black up one level and
tries again. Every other case terminates.
\item \textbf{Case 3} is a \emph{setup} move for Case 4: it rotates a red nephew
to the far side.
\item \textbf{Case 4} is the \emph{terminating} fix: a red nephew on the far side
lets one rotation pay off the debt for good.
\end{itemize}
Only Case 2 repeats, and it climbs one level each time, so the loop is
$O(\log n)$. The total number of rotations is at most \emph{three} (Case 1, then
3, then 4 in the worst chain).
\end{focusbox}

\begin{trapbox}
The notorious confusion in deletion is the ``extra black'' fiction. The node $x$
does not literally hold two colors; the doubling is a way of saying ``this path is
one black short, and we are pinning the deficit to $x$ until we find a place to
pay it.'' If you try to picture a node that is somehow black-and-a-half, you will
tie yourself in knots. Picture instead a debt token sitting on $x$, which the
fixup either hands to a red node (paid) or carries one step up the tree.
\end{trapbox}

\begin{walkbox}
On your next walk, reason about \emph{why Case 2 is the only one that loops.}
Cases 1, 3, and 4 each either end the loop or set up a strictly later case, but
Case 2 fires when both of the sibling's children are black -- the one situation
where there is no red node anywhere nearby to absorb the debt. With nothing to pay
the debt locally, the only move is to carry it up to the parent and try the whole
analysis again one level higher. Convince yourself that the debt can climb at most
to the root, where the final line of the algorithm simply forgives it.
\end{walkbox}

\begin{complexitybox}
\textbf{Deletion cost.} Finding $z$ and its successor is $O(\log n)$. The fixup
loop iterates only via Case 2, which ascends one level per iteration, so at most
$O(\log n)$ iterations, with at most three rotations total. Hence deletion is
$O(\log n)$ time, $O(1)$ rotations -- matching insertion.
\end{complexitybox}

\section{Putting It All Together}

Step back and look at what one bit per node bought us.

\begin{complexitybox}
\textbf{Red-black tree operations, all worst case.}
\begin{center}
\begin{tabular}{lcc}
\toprule
Operation & Time & Rotations \\
\midrule
\textsc{Search} / \textsc{Minimum} / \textsc{Maximum} & $O(\log n)$ & $0$ \\
\textsc{Successor} / \textsc{Predecessor} & $O(\log n)$ & $0$ \\
\textsc{Insert} & $O(\log n)$ & $\le 2$ \\
\textsc{Delete} & $O(\log n)$ & $\le 3$ \\
\bottomrule
\end{tabular}
\end{center}
The height is at most $2\log_2(n+1)$, so nothing is ever worse than logarithmic --
no matter the order of operations. There is no bad input.
\end{complexitybox}

\begin{aestheticbox}
Red-black trees are the workhorse balanced BST of the standard library era, but
they are not the only way to win this game. AVL trees keep an even tighter height
bound (about $1.44\log_2 n$) and so search a touch faster, at the cost of doing
more rotations on update -- a good trade when you read far more than you write.
B-trees (the subject of a later chapter) generalize the idea to fat nodes that hold many keys, which
is exactly what you want when each node access is a slow disk read. Skip lists
reach the same guarantees with randomization and simpler code. The red-black tree
sits at a sweet spot: a strict worst-case guarantee, few rotations per update, and
a structure that maps neatly onto a single color bit you can often hide inside an
already-present pointer.
\end{aestheticbox}

\begin{historybox}
Red-black trees were introduced by Leonidas Guibas and Robert Sedgewick in 1978
as a ``symmetric binary B-tree'' reframed with colors; the red/black coloring made
the invariants and case analysis dramatically easier to state and verify. They are
in fact a 1-to-1 disguise for 2-3-4 trees (a B-tree of order 4): a black node
together with its red children corresponds to one fat B-tree node holding up to
three keys, which is the cleanest way to \emph{understand} where the five
properties come from. For the rigorous treatment, the precise pseudocode this
chapter is built on, and the deletion proof in full, see CLRS Chapter 13
(``Red-Black Trees''), with the binary-search-tree groundwork in Chapter 12 and
the 2-3-4 / B-tree connection in Chapter 18. The original idea traces back through
Rudolf Bayer's 1972 symmetric binary B-trees.
\end{historybox}

The correspondence is worth seeing once. A black node with one or two red children
collapses into a single 2-3-4 node holding two or three keys; a lone black node is a
2-node holding one key. Property 5 (equal black-height) is just the statement that a
B-tree keeps all its leaves on the same level, and property 4 (no two reds in a row)
is what stops a B-tree node from overflowing past three keys.

\begin{center}
\begin{tikzpicture}[
  black/.style={circle,draw,fill=black,text=white,minimum size=6mm,inner sep=0pt,font=\scriptsize},
  red/.style={circle,draw=trapred,very thick,fill=red!12,minimum size=6mm,inner sep=0pt,font=\scriptsize},
  bnode/.style={rectangle,draw,rounded corners=1pt,fill=structurepurple!12,draw=structurepurple,minimum height=6mm,inner sep=3pt,font=\scriptsize},
  every edge/.style={draw,-{Latex[length=1.2mm]}}]
% red-black cluster: black A with red children a and c
\node[red]   (a) at (-1,-0.0) {$a$};
\node[black] (m) at (0,0.9) {$b$};
\node[red]   (c) at (1,0.0) {$c$};
\draw (m) edge (a); \draw (m) edge (c);
\node[font=\scriptsize] at (2.6,0.45) {$\equiv$};
% equivalent 2-3-4 node
\node[bnode] (B) at (4.6,0.45) {$a\ \,b\ \,c$};
\node[font=\scriptsize,align=center] at (0,-1.0) {one black node\\ with two red children};
\node[font=\scriptsize,align=center] at (4.6,-1.0) {one 2-3-4 node\\ holding three keys};
\end{tikzpicture}
\end{center}

\begin{focusbox}
\textbf{What to hold onto from this chapter.}
\begin{enumerate}
\item Balance is everything: BST operations cost $\Th(\text{height})$, so a tree
that stays $O(\log n)$ tall is fast forever.
\item Five properties enforce balance, and the two that matter are ``no two reds
in a row'' (property 4) and ``equal black-height on all paths'' (property 5).
Together they force the longest path to be at most twice the shortest.
\item Rotations are $O(1)$ local restructurings that preserve search order;
they are the only tool we have for changing shape.
\item Insert red, then fix up the at-most-one red-red violation with three cases
and $\le 2$ rotations.
\item Delete, and if a black node left, discharge the ``extra black'' with four
cases and $\le 3$ rotations.
\end{enumerate}
\end{focusbox}

\section*{Exercises}
\addcontentsline{toc}{section}{Exercises}

\begin{exercise}[Warmup: legal or not]\label{ch13:ex-legal}
For each of the following, decide whether it can be a valid red-black tree and
justify with the specific property that fails if not.
(a) A tree whose root is red.
(b) A black root with two red children, each of which has two black children.
(c) A path from the root that reads black, red, red, black going downward.
\end{exercise}

\begin{exercise}[Black-height counting]\label{ch13:ex-bh}
In the worked example tree at the end of Section~\ref{ch13:sec-worked-insert} (root $20$), compute the
black-height $\mathrm{bh}(x)$ for every node $x$, remembering to count the
sentinel leaves. Verify property 5 holds at the root by listing the black count of
three distinct root-to-leaf paths.
\end{exercise}

\begin{exercise}[Rotations preserve order]\label{ch13:ex-rotate}
Take the tree from the rotation lab in Section~\ref{ch13:sec-rotations}. Perform a left rotation
around the root, then a right rotation around the new root. You should recover the
original tree. Write out the in-order traversal after each step to confirm the key
order never changes, and explain in one sentence why \textsc{Right-Rotate} is the
exact inverse of \textsc{Left-Rotate}.
\end{exercise}

\begin{exercise}[Insert and trace]\label{ch13:ex-insert}
Starting from an empty tree, insert the keys $41, 38, 31, 12, 19, 8$ one at a
time. For each insertion, name which case of \textsc{RB-Insert-Fixup} fires (if
any) and draw the resulting tree with colors. (Several of these will trigger
rotations -- this sequence is designed to.)
\end{exercise}

\begin{exercise}[Why red, not black?]\label{ch13:ex-why-red}
Suppose you modified \textsc{RB-Insert} to color the new node \emph{black} instead
of red. Which of the five properties would be violated immediately upon insertion
into a non-empty tree, and why would the repair be more involved than in the red
case? Argue in terms of black-height.
\end{exercise}

\begin{exercise}[Maximum and minimum reds]\label{ch13:ex-red-count}
Among all valid red-black trees on $n = 7$ internal nodes, what is the largest
possible number of red nodes? The smallest? Sketch a tree achieving each extreme.
(Hint: think about property 4 for the maximum and property 5 for the minimum.)
\end{exercise}

\begin{exercise}[Delete and trace]\label{ch13:ex-delete}
Build the red-black tree by inserting $5, 10, 15, 20, 25, 30, 35$ (this gives a
nicely balanced tree). Then delete $10$. Identify the spliced-out node $y$, its
original color, the node $x$ that moves up, and which deletion-fixup case (if any)
runs. Draw the final tree.
\end{exercise}

\begin{exercise}[The extra-black invariant]\label{ch13:ex-extra-black}
Explain, in your own words and in at most six sentences, the loop invariant
maintained by \textsc{RB-Delete-Fixup}: what is true about the tree at the top of
each iteration, in what sense is property 5 ``violated only at $x$,'' and why does
terminating the loop (either reaching a red $x$ or the root) restore all five
properties when we finally paint $x$ black?
\end{exercise}

\begin{exercise}[AVL versus red-black, conceptually]\label{ch13:ex-avl}
An AVL tree guarantees height at most about $1.44\log_2 n$, tighter than the
red-black bound of $2\log_2(n+1)$. Without looking anything up, argue from the
balance conditions why an AVL tree should generally perform \emph{more} rotations
per insertion/deletion than a red-black tree, and describe one workload where each
structure would be the better choice.
\end{exercise}

\begin{exercise}[Challenge: join two trees]\label{ch13:ex-join}
Suppose you have two red-black trees $T_1$ and $T_2$ together with a key $k$ such
that every key in $T_1$ is less than $k$ and every key in $T_2$ is greater than
$k$. Describe an $O(\log n)$ algorithm to combine them into a single red-black tree
containing all the keys of $T_1$, all of $T_2$, and $k$. (Hint: the black-heights
of the two trees are the lever. Walk down the taller tree's appropriate spine until
you reach a node whose black-height matches the shorter tree, then splice in $k$ as
a red node joining the two and run a single insert-style fixup.)
\end{exercise}
